\documentclass[10pt,a4paper]{amsart}
\usepackage[margin=19mm]{geometry}
\usepackage{amsmath,amssymb,mathtools}
\usepackage[scr=rsfs,cal=euler]{mathalfa}
\usepackage{xfrac}
\usepackage[unicode,hidelinks,bookmarks=false]{hyperref}
\newcommand{\ie}{i.e.\ }
\newcommand{\cf}{cf.\ }
\newcommand{\resp}{respectively\ }
\newcommand{\Subsection}{\subsection}
\providecommand{\nobreakdash}{\nobreak\hbox{-}}
\setlength{\emergencystretch}{2em}
\multlinegap0pt
\usepackage{amssymb}
\usepackage{mathtools}
\usepackage{xcolor}
\newtheorem{theorem}{Theorem}
\newtheorem{corollary}{Corollary}
\newcommand{\BH}{\mathsf B}
\title{Polynomial growth of Bohnenblust--Hille constants on the Hamming cube}
\author{Paata Ivanisvili}
\date{Source: arXiv:2609.12427v1 (CC BY 4.0)}
\begin{document}
\maketitle

\noindent\emph{Source and licence.} Polynomial growth of Bohnenblust--Hille constants on the Hamming cube. Version arXiv:2609.12427v1 (CC BY 4.0), distributed by arXiv under the Creative Commons Attribution 4.0 International licence, http://creativecommons.org/licenses/by/4.0/, as recorded in the arXiv licence metadata for this exact version and shown on the abstract page https://arxiv.org/abs/2609.12427v1. Full licence notice: \texttt{licenses/arxiv-2609.12427-CC-BY-4.0.txt}.\par\smallskip

\noindent\textit{Editorial scope.} Counted unit: the article's corollary cor:AA (source0.tex lines 154--163). The article's complete original proof of it is delivered with it (source0.tex lines 662--687, one \texttt{proof} environment of 26 lines, which the article places away from the statement). No mathematics is written, repaired or re-derived: the statement and the proof are the article's own lines, in the article's own notation, and no retained line is substituted or re-flowed.  Retained uncounted as the source-local context the proof needs: lines 114--123; lines 131--138.  Nothing was omitted from any retained span, so the package carries no omission record.  No retained line contains a citation, so the wrapper carries no bibliography and no bibliography entry.  No retained line contains a figure environment, an image-inclusion command or drawing code, so no image is delivered and the package declares no asset dependency.  Editorial formatting only, in addition: an AMS \texttt{amsart} preamble that restates the article's own declarations and macros used by the retained lines, namely the environments \texttt{theorem}, \texttt{corollary} on separate counters, together with the standard packages those lines need.  The retained environments print in this document as Theorem 1, Corollary 1 in order of appearance, whereas the article prints the same items with the numbering of its own full document.  The complete native source of the article as distributed in the arXiv e-print for this exact version is bundled unchanged as \texttt{sources/210115/source0.tex}.
\begin{theorem}\label{thm:main}
There is an absolute constant $K<\infty$ such that, for every pair of
positive integers $m,n$ with $m\le n$ and every $f:Q_n\to\mathbb C$ of degree
at most $m$,
\begin{equation}\label{eq:main}
\left(\sum_{|S|\le m}|\widehat f(S)|^{2m/(m+1)}\right)^{(m+1)/(2m)}
\le Km^{27}\|f\|_\infty.
\end{equation}
Equivalently, $\BH_m\le Km^{27}$ for every $m\ge1$.
\end{theorem}

\begin{theorem}\label{thm:weighted}
There is an absolute constant $K<\infty$ such that every polynomial as in
Theorem~\ref{thm:main} satisfies
\begin{equation}\label{eq:weighted-intro}
\left(\sum_{r=1}^m\frac{\mathcal{S}_{r}(f)^2}{r^{10}}\right)^{1/2}
\le Km^{22}\|f\|_\infty.
\end{equation}
\end{theorem}

\begin{corollary}\label{cor:AA}
There is an absolute constant $c>0$ such that, for every $1\le m\le n$
and every $f:Q_n\to[-1,1]$ of degree at most $m$ satisfying
$|\widehat f(S)|=a_{|S|}$ for some $a_0,\ldots,a_m\ge0$ and all $|S|\le m$,
\[
\operatorname{Inf}_i(f)\ge
\frac{c}{m^{54}}\operatorname{Var}(f)^2
\qquad\text{for every }i\in[n].
\]
\end{corollary}

\begin{proof}[Proof of Corollary~\ref{cor:AA}]
Put $v_r=\sum_{|S|=r}|\widehat f(S)|^2=\binom nr a_r^2$. The hypothesis
gives, for every $i\in[n]$,
\[
\operatorname{Var}(f)=\sum_{r=1}^m v_r,
\qquad
\operatorname{Inf}_i(f)=\sum_{r=1}^m\binom{n-1}{r-1}a_r^2
=\frac1n\sum_{r=1}^m rv_r.
\]
Moreover, since $2/q_r=1+1/r$ and $\binom nr\ge(n/r)^r$,
\[
\mathcal{S}_r(f)^2=\binom nr^{1/r}v_r\ge\frac nr\,v_r.
\]
Cauchy--Schwarz and Theorem~\ref{thm:weighted} therefore yield
\begin{align*}
\operatorname{Var}(f)^2
&\le\left(\sum_{r=1}^m rv_r\right)
       \left(\sum_{r=1}^m\frac{v_r}{r}\right)\\
&\le\operatorname{Inf}_i(f)\sum_{r=1}^m\mathcal{S}_r(f)^2\\
&\le m^{10}\operatorname{Inf}_i(f)
       \sum_{r=1}^m\frac{\mathcal{S}_r(f)^2}{r^{10}}\\
&\le K^2m^{54}\operatorname{Inf}_i(f),
\end{align*}
where the last step uses $\|f\|_\infty\le1$.
Taking $c=K^{-2}$ proves the corollary.
\end{proof}

\end{document}
